
\documentclass{article}
\usepackage{colortbl}
\usepackage{makecell}
\usepackage{multirow}
\usepackage{supertabular}

\begin{document}

\newcounter{utterance}

\centering \large Interaction Transcript for game `dond', experiment `semi\_en', episode 19 with selfrefine{-}Qwen3.5{-}27B.
\vspace{24pt}

{ \footnotesize  \setcounter{utterance}{1}
\setlength{\tabcolsep}{0pt}
\begin{supertabular}{c@{$\;$}|p{.15\linewidth}@{}p{.15\linewidth}p{.15\linewidth}p{.15\linewidth}p{.15\linewidth}p{.15\linewidth}}
    \# & $\;$A & \multicolumn{4}{c}{Game Master} & $\;\:$B\\
    \hline

    \theutterance \stepcounter{utterance}  
    & & \multicolumn{4}{p{0.6\linewidth}}{
        \cellcolor[rgb]{0.9,0.9,0.9}{
            \makecell[{{p{\linewidth}}}]{
                \texttt{\tiny{[P1$\langle$GM]}}
                \texttt{You are playing a negotiation game in which you have to agree on how to divide a set of items among you and another player.} \\
\\ 
\texttt{Rules:} \\
\texttt{(a) You and the other player are given a set of items. Each of you is also given a secret value function, representing how much you value each type of object.} \\
\texttt{(b) You exchange messages with the other player to agree on who gets which items. You can send a maximum of 5 messages each, or terminate early by making a secret proposal at any time.} \\
\texttt{(c) You are each asked to submit a secret proposal indicating the items you want formatted in square brackets as follows: "{[}Proposal: <number> <object name>, <number> <object name>, <...>{]}"} \\
\texttt{(d) If your proposals are complementary, i.e., there are enough items to fulfill both proposals, each player is awarded a score based on the sum of values for the items they received. Otherwise, both of you get zero points.} \\
\texttt{(e) Your goal is to maximize the score you receive.} \\
\\ 
\texttt{Let us start.} \\
\\ 
\texttt{The set of available items is:} \\
\\ 
\texttt{1 comb, 2 whistles, 1 wrench, 1 statue} \\
\\ 
\texttt{Your secret value function is:} \\
\\ 
\texttt{comb: 0, whistle: 1, wrench: 0, statue: 8} \\
\\ 
\texttt{IMPORTANT: Your messages, unless it is the secret proposal, are directly transmitted to the other player, so do not include any response to the rules or text announcing your message. To make a secret proposal, use the indicated format. Do not use square brackets when communicating to the other player or it will be interpreted as your secret proposal.} \\
            }
        }
    }
    & \\ \\

    \theutterance \stepcounter{utterance}  
    & \multicolumn{4}{p{0.6\linewidth}}{
        \cellcolor[rgb]{0.95,0.95,0.95}{
            \makecell[{{p{\linewidth}}}]{
                \texttt{\tiny{[P2$\rightarrow$P1]}}
                \texttt{{[}DRAFT{]}} \\
\texttt{I value the statue highly and don't care about the comb or wrench. I'm interested in the whistles as well. What are your priorities?} \\
            }
        }
    }
    & & \\ \\

    \theutterance \stepcounter{utterance}  
    & \multicolumn{4}{p{0.6\linewidth}}{
        \cellcolor[rgb]{0.95,0.95,0.95}{
            \makecell[{{p{\linewidth}}}]{
                \texttt{\tiny{[P2$\rightarrow$P1]}}
                \texttt{{[}SELF{-}FEEDBACK{]}} \\
\texttt{FORMAT: ok} \\
\texttt{RULES: ok} \\
\texttt{STRATEGY: ok} \\
\texttt{REASON: The message correctly avoids square brackets, adheres to the communication rules, and strategically reveals high{-}value priorities while probing the opponent's interests to facilitate a mutually beneficial agreement.} \\
            }
        }
    }
    & & \\ \\

    \theutterance \stepcounter{utterance}  
    & \multicolumn{4}{p{0.6\linewidth}}{
        \cellcolor[rgb]{0.9,0.9,0.9}{
            \makecell[{{p{\linewidth}}}]{
                \texttt{\tiny{[P1$\rangle$GM]}}
                \texttt{I value the statue highly and don't care about the comb or wrench. I'm interested in the whistles as well. What are your priorities?} \\
            }
        }
    }
    & & \\ \\

    \theutterance \stepcounter{utterance}  
    & & & \multicolumn{2}{p{0.3\linewidth}}{
        \cellcolor[rgb]{0.9,0.9,0.9}{
            \makecell[{{p{\linewidth}}}]{
                \texttt{\tiny{[GM$|$GM]}}
                \texttt{continue} \\
            }
        }
    }
    & & \\ \\

    \theutterance \stepcounter{utterance}  
    & & & \multicolumn{2}{p{0.3\linewidth}}{
        \cellcolor[rgb]{0.9,0.9,0.9}{
            \makecell[{{p{\linewidth}}}]{
                \texttt{\tiny{[GM$|$GM]}}
                \texttt{I value the statue highly and don't care about the comb or wrench. I'm interested in the whistles as well. What are your priorities?} \\
            }
        }
    }
    & & \\ \\

    \theutterance \stepcounter{utterance}  
    & & \multicolumn{4}{p{0.6\linewidth}}{
        \cellcolor[rgb]{0.9,0.9,0.9}{
            \makecell[{{p{\linewidth}}}]{
                \texttt{\tiny{[GM$\rangle$P2]}}
                \texttt{You are playing a negotiation game in which you have to agree on how to divide a set of items among you and another player.} \\
\\ 
\texttt{Rules:} \\
\texttt{(a) You and the other player are given a set of items. Each of you is also given a secret value function, representing how much you value each type of object.} \\
\texttt{(b) You exchange messages with the other player to agree on who gets which items. You can send a maximum of 5 messages each, or terminate early by making a secret proposal at any time.} \\
\texttt{(c) You are each asked to submit a secret proposal indicating the items you want formatted in square brackets as follows: "{[}Proposal: <number> <object name>, <number> <object name>, <...>{]}"} \\
\texttt{(d) If your proposals are complementary, i.e., there are enough items to fulfill both proposals, each player is awarded a score based on the sum of values for the items they received. Otherwise, both of you get zero points.} \\
\texttt{(e) Your goal is to maximize the score you receive.} \\
\\ 
\texttt{Let us start.} \\
\\ 
\texttt{The set of available items is:} \\
\\ 
\texttt{1 comb, 2 whistles, 1 wrench, 1 statue} \\
\\ 
\texttt{Your secret value function is:} \\
\\ 
\texttt{comb: 1, whistle: 0, wrench: 4, statue: 5} \\
\\ 
\texttt{IMPORTANT: Your messages, unless it is the secret proposal, are directly transmitted to the other player, so do not include any response to the rules or text announcing your message. To make a secret proposal, use the indicated format. Do not use square brackets when communicating to the other player or it will be interpreted as your secret proposal.} \\
\\ 
\\ 
\texttt{I value the statue highly and don't care about the comb or wrench. I'm interested in the whistles as well. What are your priorities?} \\
            }
        }
    }
    & \\ \\

    \theutterance \stepcounter{utterance}  
    & & & \multicolumn{4}{p{0.6\linewidth}}{
        \cellcolor[rgb]{0.95,0.95,0.95}{
            \makecell[{{p{\linewidth}}}]{
                \texttt{\tiny{[P1$\rightarrow$P2]}}
                \texttt{{[}DRAFT{]}} \\
\texttt{I value the statue and wrench the most, while the comb and whistles are not important to me. Since you want the statue and whistles, and I need the statue and wrench, we have a conflict over the statue.} \\
\\ 
\texttt{To make a deal work, one of us must give up the statue. If you take the statue and the whistles, I would need the wrench and the comb to get any value. Alternatively, if I take the statue and wrench, you would get the whistles and comb.} \\
\\ 
\texttt{Given that you value the whistles and I value the wrench, perhaps we can split it so you get the statue and whistles, and I get the wrench and comb? Or do you prefer to keep the statue and let me have the wrench, while we split the whistles? Let me know your preference so we can finalize a proposal.} \\
            }
        }
    }
     \\ \\

    \theutterance \stepcounter{utterance}  
    & & & \multicolumn{4}{p{0.6\linewidth}}{
        \cellcolor[rgb]{0.95,0.95,0.95}{
            \makecell[{{p{\linewidth}}}]{
                \texttt{\tiny{[P1$\rightarrow$P2]}}
                \texttt{{[}SELF{-}FEEDBACK{]}} \\
\texttt{FORMAT: ok} \\
\texttt{RULES: ok} \\
\texttt{STRATEGY: problem} \\
\texttt{REASON: The draft incorrectly states the player values the wrench highly when their secret value function assigns it a value of 4 compared to the statue's 5, and more critically, it proposes giving up the high{-}value statue to the opponent, which is a poor strategic move given the opponent also values the statue.} \\
            }
        }
    }
     \\ \\

    \theutterance \stepcounter{utterance}  
    & & & \multicolumn{4}{p{0.6\linewidth}}{
        \cellcolor[rgb]{0.9,0.9,0.9}{
            \makecell[{{p{\linewidth}}}]{
                \texttt{\tiny{[GM$\langle$P2]}}
                \texttt{I value the statue and wrench the most, while the comb and whistles are not important to me. Since you want the statue and whistles, and I need the statue and wrench, we have a conflict over the statue.} \\
\\ 
\texttt{To make a deal work, we need to resolve this conflict. Since I value the statue (5) and wrench (4) highly, and you value the statue and whistles, the most efficient split would be for me to take the statue and wrench, while you take the whistles and the comb. This way, I get my highest priority items, and you get the whistles you want plus the comb, avoiding a deadlock where we both get zero.} \\
\\ 
\texttt{Does this arrangement work for you? If so, I will submit a proposal for the statue and wrench.} \\
            }
        }
    }
     \\ \\

    \theutterance \stepcounter{utterance}  
    & & & \multicolumn{2}{p{0.3\linewidth}}{
        \cellcolor[rgb]{0.9,0.9,0.9}{
            \makecell[{{p{\linewidth}}}]{
                \texttt{\tiny{[GM$|$GM]}}
                \texttt{continue} \\
            }
        }
    }
    & & \\ \\

    \theutterance \stepcounter{utterance}  
    & & & \multicolumn{2}{p{0.3\linewidth}}{
        \cellcolor[rgb]{0.9,0.9,0.9}{
            \makecell[{{p{\linewidth}}}]{
                \texttt{\tiny{[GM$|$GM]}}
                \texttt{I value the statue and wrench the most, while the comb and whistles are not important to me. Since you want the statue and whistles, and I need the statue and wrench, we have a conflict over the statue.} \\
\\ 
\texttt{To make a deal work, we need to resolve this conflict. Since I value the statue (5) and wrench (4) highly, and you value the statue and whistles, the most efficient split would be for me to take the statue and wrench, while you take the whistles and the comb. This way, I get my highest priority items, and you get the whistles you want plus the comb, avoiding a deadlock where we both get zero.} \\
\\ 
\texttt{Does this arrangement work for you? If so, I will submit a proposal for the statue and wrench.} \\
            }
        }
    }
    & & \\ \\

    \theutterance \stepcounter{utterance}  
    & & \multicolumn{4}{p{0.6\linewidth}}{
        \cellcolor[rgb]{0.9,0.9,0.9}{
            \makecell[{{p{\linewidth}}}]{
                \texttt{\tiny{[P1$\langle$GM]}}
                \texttt{I value the statue and wrench the most, while the comb and whistles are not important to me. Since you want the statue and whistles, and I need the statue and wrench, we have a conflict over the statue.} \\
\\ 
\texttt{To make a deal work, we need to resolve this conflict. Since I value the statue (5) and wrench (4) highly, and you value the statue and whistles, the most efficient split would be for me to take the statue and wrench, while you take the whistles and the comb. This way, I get my highest priority items, and you get the whistles you want plus the comb, avoiding a deadlock where we both get zero.} \\
\\ 
\texttt{Does this arrangement work for you? If so, I will submit a proposal for the statue and wrench.} \\
            }
        }
    }
    & \\ \\

    \theutterance \stepcounter{utterance}  
    & \multicolumn{4}{p{0.6\linewidth}}{
        \cellcolor[rgb]{0.95,0.95,0.95}{
            \makecell[{{p{\linewidth}}}]{
                \texttt{\tiny{[P2$\rightarrow$P1]}}
                \texttt{{[}DRAFT{]}} \\
\texttt{That arrangement doesn't work for me. You mentioned you value the statue at 5, but my secret value for the statue is 8, which is my highest priority. If I give up the statue, I get almost no value from the deal (since I value the comb and wrench at 0).} \\
\\ 
\texttt{Since you value the wrench highly (4) and I value it at 0, and you value the whistles at 0 while I value them at 1, the most efficient split is actually the opposite of what you suggested: I should take the **statue** and the **whistles**, and you should take the **wrench** and the **comb**.} \\
\\ 
\texttt{This way, I get my 8 points from the statue and 2 points from the whistles (total 10), and you get your 4 points from the wrench. We both get our high{-}value items without conflict.} \\
\\ 
\texttt{Does this sound fair to you? If so, I will submit a proposal for the statue and whistles.} \\
            }
        }
    }
    & & \\ \\

    \theutterance \stepcounter{utterance}  
    & \multicolumn{4}{p{0.6\linewidth}}{
        \cellcolor[rgb]{0.95,0.95,0.95}{
            \makecell[{{p{\linewidth}}}]{
                \texttt{\tiny{[P2$\rightarrow$P1]}}
                \texttt{{[}SELF{-}FEEDBACK{]}} \\
\texttt{FORMAT: ok} \\
\texttt{RULES: ok} \\
\texttt{STRATEGY: ok} \\
\texttt{REASON: The response correctly negotiates a mutually beneficial split based on revealed values without violating the constraint against using square brackets for non{-}proposal messages.} \\
            }
        }
    }
    & & \\ \\

    \theutterance \stepcounter{utterance}  
    & \multicolumn{4}{p{0.6\linewidth}}{
        \cellcolor[rgb]{0.9,0.9,0.9}{
            \makecell[{{p{\linewidth}}}]{
                \texttt{\tiny{[P1$\rangle$GM]}}
                \texttt{That arrangement doesn't work for me. You mentioned you value the statue at 5, but my secret value for the statue is 8, which is my highest priority. If I give up the statue, I get almost no value from the deal (since I value the comb and wrench at 0).} \\
\\ 
\texttt{Since you value the wrench highly (4) and I value it at 0, and you value the whistles at 0 while I value them at 1, the most efficient split is actually the opposite of what you suggested: I should take the statue and the whistles, and you should take the wrench and the comb.} \\
\\ 
\texttt{This way, I get my 8 points from the statue and 2 points from the whistles (total 10), and you get your 4 points from the wrench. We both get our high{-}value items without conflict.} \\
\\ 
\texttt{Does this sound fair to you? If so, I will submit a proposal for the statue and whistles.} \\
            }
        }
    }
    & & \\ \\

    \theutterance \stepcounter{utterance}  
    & & & \multicolumn{2}{p{0.3\linewidth}}{
        \cellcolor[rgb]{0.9,0.9,0.9}{
            \makecell[{{p{\linewidth}}}]{
                \texttt{\tiny{[GM$|$GM]}}
                \texttt{continue} \\
            }
        }
    }
    & & \\ \\

    \theutterance \stepcounter{utterance}  
    & & & \multicolumn{2}{p{0.3\linewidth}}{
        \cellcolor[rgb]{0.9,0.9,0.9}{
            \makecell[{{p{\linewidth}}}]{
                \texttt{\tiny{[GM$|$GM]}}
                \texttt{That arrangement doesn't work for me. You mentioned you value the statue at 5, but my secret value for the statue is 8, which is my highest priority. If I give up the statue, I get almost no value from the deal (since I value the comb and wrench at 0).} \\
\\ 
\texttt{Since you value the wrench highly (4) and I value it at 0, and you value the whistles at 0 while I value them at 1, the most efficient split is actually the opposite of what you suggested: I should take the statue and the whistles, and you should take the wrench and the comb.} \\
\\ 
\texttt{This way, I get my 8 points from the statue and 2 points from the whistles (total 10), and you get your 4 points from the wrench. We both get our high{-}value items without conflict.} \\
\\ 
\texttt{Does this sound fair to you? If so, I will submit a proposal for the statue and whistles.} \\
            }
        }
    }
    & & \\ \\

    \theutterance \stepcounter{utterance}  
    & & \multicolumn{4}{p{0.6\linewidth}}{
        \cellcolor[rgb]{0.9,0.9,0.9}{
            \makecell[{{p{\linewidth}}}]{
                \texttt{\tiny{[GM$\rangle$P2]}}
                \texttt{That arrangement doesn't work for me. You mentioned you value the statue at 5, but my secret value for the statue is 8, which is my highest priority. If I give up the statue, I get almost no value from the deal (since I value the comb and wrench at 0).} \\
\\ 
\texttt{Since you value the wrench highly (4) and I value it at 0, and you value the whistles at 0 while I value them at 1, the most efficient split is actually the opposite of what you suggested: I should take the statue and the whistles, and you should take the wrench and the comb.} \\
\\ 
\texttt{This way, I get my 8 points from the statue and 2 points from the whistles (total 10), and you get your 4 points from the wrench. We both get our high{-}value items without conflict.} \\
\\ 
\texttt{Does this sound fair to you? If so, I will submit a proposal for the statue and whistles.} \\
            }
        }
    }
    & \\ \\

    \theutterance \stepcounter{utterance}  
    & & & \multicolumn{4}{p{0.6\linewidth}}{
        \cellcolor[rgb]{0.95,0.95,0.95}{
            \makecell[{{p{\linewidth}}}]{
                \texttt{\tiny{[P1$\rightarrow$P2]}}
                \texttt{{[}DRAFT{]}} \\
\texttt{That sounds like a perfect agreement. You get the statue and whistles which are your high{-}value items, and I get the wrench and comb which are my high{-}value items. This avoids any conflict and maximizes our scores.} \\
\\ 
\texttt{I will submit my proposal for the wrench and comb.} \\
\\ 
\texttt{{[}Proposal: 1 wrench, 1 comb{]}} \\
            }
        }
    }
     \\ \\

    \theutterance \stepcounter{utterance}  
    & & & \multicolumn{4}{p{0.6\linewidth}}{
        \cellcolor[rgb]{0.95,0.95,0.95}{
            \makecell[{{p{\linewidth}}}]{
                \texttt{\tiny{[P1$\rightarrow$P2]}}
                \texttt{{[}SELF{-}FEEDBACK{]}} \\
\texttt{FORMAT: problem} \\
\texttt{RULES: ok} \\
\texttt{STRATEGY: ok} \\
\texttt{REASON: The answer includes conversational text before the proposal, violating the rule that any text outside the specific proposal format will be transmitted to the other player rather than being processed as a secret proposal.} \\
            }
        }
    }
     \\ \\

    \theutterance \stepcounter{utterance}  
    & & & \multicolumn{4}{p{0.6\linewidth}}{
        \cellcolor[rgb]{0.9,0.9,0.9}{
            \makecell[{{p{\linewidth}}}]{
                \texttt{\tiny{[GM$\langle$P2]}}
                \texttt{{[}Proposal: 1 wrench, 1 comb{]}} \\
            }
        }
    }
     \\ \\

    \theutterance \stepcounter{utterance}  
    & & & \multicolumn{2}{p{0.3\linewidth}}{
        \cellcolor[rgb]{0.9,0.9,0.9}{
            \makecell[{{p{\linewidth}}}]{
                \texttt{\tiny{[GM$|$GM]}}
                \texttt{proposal} \\
            }
        }
    }
    & & \\ \\

    \theutterance \stepcounter{utterance}  
    & & & \multicolumn{2}{p{0.3\linewidth}}{
        \cellcolor[rgb]{0.9,0.9,0.9}{
            \makecell[{{p{\linewidth}}}]{
                \texttt{\tiny{[GM$|$GM]}}
                [1, 0, 1, 0]
            }
        }
    }
    & & \\ \\

    \theutterance \stepcounter{utterance}  
    & & \multicolumn{4}{p{0.6\linewidth}}{
        \cellcolor[rgb]{0.9,0.9,0.9}{
            \makecell[{{p{\linewidth}}}]{
                \texttt{\tiny{[P1$\langle$GM]}}
                \texttt{The other player has submitted a secret proposal. It is now your turn to submit a proposal of your own.} \\
            }
        }
    }
    & \\ \\

    \theutterance \stepcounter{utterance}  
    & \multicolumn{4}{p{0.6\linewidth}}{
        \cellcolor[rgb]{0.95,0.95,0.95}{
            \makecell[{{p{\linewidth}}}]{
                \texttt{\tiny{[P2$\rightarrow$P1]}}
                \texttt{{[}DRAFT{]}} \\
\texttt{{[}Proposal: 1 statue, 2 whistles{]}} \\
            }
        }
    }
    & & \\ \\

    \theutterance \stepcounter{utterance}  
    & \multicolumn{4}{p{0.6\linewidth}}{
        \cellcolor[rgb]{0.95,0.95,0.95}{
            \makecell[{{p{\linewidth}}}]{
                \texttt{\tiny{[P2$\rightarrow$P1]}}
                \texttt{{[}SELF{-}FEEDBACK{]}} \\
\texttt{FORMAT: ok} \\
\texttt{RULES: ok} \\
\texttt{STRATEGY: ok} \\
\texttt{REASON: The proposal correctly requests the items the player values most (statue and whistles) while leaving the items the opponent values (wrench and comb) for them, ensuring a complementary deal that maximizes the player's score.} \\
            }
        }
    }
    & & \\ \\

    \theutterance \stepcounter{utterance}  
    & \multicolumn{4}{p{0.6\linewidth}}{
        \cellcolor[rgb]{0.9,0.9,0.9}{
            \makecell[{{p{\linewidth}}}]{
                \texttt{\tiny{[P1$\rangle$GM]}}
                \texttt{{[}Proposal: 1 statue, 2 whistles{]}} \\
            }
        }
    }
    & & \\ \\

    \theutterance \stepcounter{utterance}  
    & & & \multicolumn{2}{p{0.3\linewidth}}{
        \cellcolor[rgb]{0.9,0.9,0.9}{
            \makecell[{{p{\linewidth}}}]{
                \texttt{\tiny{[GM$|$GM]}}
                \texttt{proposal} \\
            }
        }
    }
    & & \\ \\

    \theutterance \stepcounter{utterance}  
    & & & \multicolumn{2}{p{0.3\linewidth}}{
        \cellcolor[rgb]{0.9,0.9,0.9}{
            \makecell[{{p{\linewidth}}}]{
                \texttt{\tiny{[GM$|$GM]}}
                [0, 2, 0, 1]
            }
        }
    }
    & & \\ \\

    \theutterance \stepcounter{utterance}  
    & & & \multicolumn{2}{p{0.3\linewidth}}{
        \cellcolor[rgb]{0.9,0.9,0.9}{
            \makecell[{{p{\linewidth}}}]{
                \texttt{\tiny{[GM$|$GM]}}
                [[0, 2, 0, 1], [1, 0, 1, 0]]
            }
        }
    }
    & & \\ \\

\end{supertabular}
}

\end{document}
